\begin{figure*}[h]
    \begin{AIbox}{Autoformalization Task Prompt}

\begin{quote}
\ttfamily
\tiny
\vspace{2mm}

YOUR TASK --- build a function for the clinical concept `\{concept\_name\}':

You need to write a Python function that can help answer clinical
questions about `\{concept\_name\}' for individual patients.

Here are example questions your function will need to support:
\{sample\_questions\}

APPROACH:
1. **Search for guidelines** --- call `search\_guidelines()' to see what clinical
   guidelines are available, then use relevant keywords (e.g.
   `search\_guidelines("\{concept\_name\}")') to find the most relevant ones.
   Read each relevant guideline with `get\_guideline(name)' --- these contain
   clinical definitions, scoring criteria, and implementation details that
   are essential for getting the logic right.
2. **Check for reusable functions** --- call `search\_functions()' to see if any
   previously implemented functions could serve as helpers. Use
   `get\_function\_info(name)' to inspect their signatures and docstrings,
   and `load\_function(name)' to import ones you want to reuse.
3. **Explore the database schema** --- list tables, inspect columns, look up
   dictionary/lookup tables (d\_labitems, d\_items, d\_icd\_diagnoses, etc.).
   Run sample queries to understand data formats.
4. **Write a Python function** --- implement the clinical logic based on what
   you learned from the guidelines and the database schema.  The function
   should accept patient identifiers (e.g. subject\_id, hadm\_id, stay\_id)
   and return information that captures the clinical concept.
   The output format is flexible --- return whatever best represents the
   concept (e.g., a value, dict, DataFrame, boolean, list, etc.).
5. **Test thoroughly** --- call your function with several sample patients
   and verify the output looks reasonable.  Debug and fix any issues.
   Do NOT assign FINAL\_FUNCTION until you are confident it works.
6. **Save it** --- once satisfied, write a single self-contained code block
   with ALL imports, helpers, function def, and FINAL\_FUNCTION assignment.
   No test calls or print statements in that final block.

FINAL FUNCTION REQUIREMENTS:
  - It should return useful clinical information that captures the concept.
  - The input should involve some patient identifier (e.g. stay\_id)
  - The output can be any format (scalar, dict, DataFrame, etc.) --- choose
    whatever best represents the clinical concept.
  - Use `query\_db(sql)' or `query\_fhir(resource\_type, params)' to look up data.
  - Have a detailed docstring which explains the expected input and output.

Start by searching for guidelines relevant to `\{concept\_name\}'.

\end{quote}

\end{AIbox}
    \caption{The task prompt provided to the LLM agent to initiate autoformalization of a clinical concept. Template variables \{concept\_name\} and \{sample\_questions\} are filled at runtime.}
    \label{prompts:react_user_prompt}

\end{figure*}
